% TOP_PROBLEM: {"id":30007606,"problem_number":"LOCAL-30007606","title":"Net macroprudential benefits","category_id":21,"category":{"id":21,"name":"economics","display_name":"Economics","description":"Open economic theory statements, published research agendas, and proposed scoped specifications; question provenance and historical priority are qualified per record.","slug":"economics","order_index":21,"created_at":"2026-10-08T00:00:00Z"},"status":"research_question","external_url":null,"legacy_ids":[],"published":false,"statement":"In the explicitly specified setting: What are the long-run welfare, inequality and crisis-risk effects of macroprudential policies, beyond their effect on credit growth? The mathematical statement above fixes the target, assumptions and scope of this catalogue formulation.","economics":{"original_id":95,"field":"Banking, credit and financial stability","kind":"empirical","formalization_basis":"adapted_research_question","source_status":"research_agenda","priority_status":"earliest_found","priority_note":"This is the earliest explicit statement verified in this audit; first priority for the full catalogue formulation is not established.","earliest_verified_reference":{"authors":["Nina Biljanovska","Sophia Chen","Gaston Gelos","Deniz Igan","Maria Soledad Martinez Peria","Erlend Nier","Fabián Valencia"],"title":"Macroprudential Policy Effects: Evidence and Open Questions","year":2023,"url":"https://www.imf.org/en/publications/departmental-papers-policy-papers/issues/2023/03/26/macroprudential-policy-effects-evidence-and-open-questions-527293","doi":"10.5089/9798400226304.087","locator":"Overview; Core Questions; Key Findings and Themes, official publication landing-page summary","evidence_summary":"Asks leakage, resilience, activity and heterogeneous policy effects."},"current_reference":{"authors":["Nina Biljanovska","Sophia Chen","Gaston Gelos","Deniz Igan","Maria Soledad Martinez Peria","Erlend Nier","Fabián Valencia"],"title":"Macroprudential Policy Effects: Evidence and Open Questions","year":2023,"url":"https://www.imf.org/en/publications/departmental-papers-policy-papers/issues/2023/03/26/macroprudential-policy-effects-evidence-and-open-questions-527293","doi":"10.5089/9798400226304.087","locator":"Overview; Core Questions; Key Findings and Themes, official publication landing-page summary","evidence_summary":"Asks leakage, resilience, activity and heterogeneous policy effects."},"formalization_note":"The source asks resilience and activity questions; household welfare weights and full inequality incidence are additional specification choices, not directly stated source formulas."},"statusQualification":"Published question or research agenda verified at the cited locator. The mathematical specification is adapted for this catalogue and includes additional assumptions; source publication does not establish first-ever priority or certify this exact model remains unsolved. The source asks resilience and activity questions; household welfare weights and full inequality incidence are additional specification choices, not directly stated source formulas.","statusReviewedAt":"2026-10-08"}
% FIRST_POSTED: 2026-10-08
% ENTRY_KIND: statement_only
% !TeX program = lualatex
\documentclass[11pt]{article}
\usepackage{fontspec}
\usepackage[a4paper,margin=25mm]{geometry}
\usepackage{amsmath,amssymb,mathtools}
\usepackage{xurl}
\usepackage[hidelinks]{hyperref}
\newcommand{\sref}[2]{\hyperlink{source:#1:#2}{[S#2]}}
\setlength{\emergencystretch}{3em}
\begin{document}
% problemId:30007606
\section*{Net macroprudential benefits}\label{30007606}
\subsection{Definitions and mathematical statement}
Fix a target population of households with welfare weights \(\omega_i\geq0\), \(\sum_i\omega_i=1\), a discount factor \(\beta\in(0,1)\), and a feasible macroprudential policy \(a\). Household consumption under intervention \(do(a)\) is \(c_{it}(a)\); crisis losses, lending restrictions and fiscal costs enter the specified equilibrium rather than being added twice. Define
\[W(a)=E\left[\sum_{t=0}^{\infty}\beta^t\sum_i\omega_i u_i(c_{it}(a))\right],\]
assuming convergence and independently specified utility functions \(u_i\). The task is to identify or bound \(W(a)-W(0)\), its distribution across household groups, and the policy-induced change in crisis probability. Permit policy effects to depend on initial credit conditions and interact with monetary and microprudential rules. State how outcomes absent adoption are identified when regulators respond to latent risk. Account for nonbank and foreign lending, housing costs and borrowing access across income groups. A reduction in credit growth is neither a welfare benefit nor a welfare cost without these channels. Report sensitivity to welfare weights, tail-loss assumptions and discounting, and distinguish short-run causal estimates from long-run structural extrapolation.

\par\textbf{Formulation and scope.} The source asks resilience and activity questions; household welfare weights and full inequality incidence are additional specification choices, not directly stated source formulas.

\par\textbf{Open-question provenance.} An explicit question or research agenda was verified at the source locator. This does not by itself certify that every later version remains unresolved \sref{30007606}{1}.

\par\textbf{Historical priority.} This is the earliest explicit statement verified in this audit; first priority for the full catalogue formulation is not established.

\subsection{Short English statement}
In the explicitly specified setting: What are the long-run welfare, inequality and crisis-risk effects of macroprudential policies, beyond their effect on credit growth? The mathematical statement above fixes the target, assumptions and scope of this catalogue formulation.

\subsection{Sources}
\begin{itemize}\raggedright
\item[\textnormal{[S1]}]\hypertarget{source:30007606:1}{} Nina Biljanovska, Sophia Chen, Gaston Gelos, Deniz Igan, Maria Soledad Martinez Peria, Erlend Nier, Fabián Valencia. 2023. Macroprudential Policy Effects: Evidence and Open Questions. Overview; Core Questions; Key Findings and Themes, official publication landing-page summary.
\url{https://www.imf.org/en/publications/departmental-papers-policy-papers/issues/2023/03/26/macroprudential-policy-effects-evidence-and-open-questions-527293}.
DOI: 10.5089/9798400226304.087.
{\small\textit{Earliest verified question/agenda reference; historical priority qualified below. Asks leakage, resilience, activity and heterogeneous policy effects.}}
\end{itemize}
\end{document}
